\uuid{aFh9}
\titre{Domaine de définition}
\auteur{bodin ; hanani ; mzari}
\organisation{exo7}
\licence{CC BY-NC-SA 4.0}
\datecreate{2020-01-20}
\niveau{L1}
\module{Analyse}
\chapitre{Continuité, limite et étude de fonctions réelles}
\sousChapitre{Autre}
\difficulte{3}
\theme{Notion de fonctions, domaine de définition}
\exo7qcmid{lille-1.continuite.010}

\contenu{
\begin{question}
\qclassification{123.00}
\qtags{motcle=domaine de définition}
Soient $f : ]-\infty,0[ \to ]0,1[$ et $g : ]-2,2[ \to ]0,+\infty[$.
Quelles sont les assertions vraies ?
\begin{answers}
    \bad{Le domaine de définition de $x \mapsto g\big(f(2x)\big)$ est $]-1,1[$.}

    \bad{Le domaine de définition de $x \mapsto g\big( \ln (f(x)) \big)$ est $]0,+\infty[$.}

    \good{Le domaine de définition de $x \mapsto \frac{g(x+1)}{f(x)}$ est $]-3,0[$.}

    \good{Le domaine de définition de $x \mapsto \frac{f(x) \times g(x)}{f(x)+g(x)}$ est $]-2,0[$.}
\end{answers}
\begin{explanations}
Si $x \mapsto f(x)$ est définie sur $]a,b[$ alors
$x \mapsto f(x+k)$ est définie sur $]a-k,b-k[$
et $x \mapsto f(\ell x)$ est définie sur $]\frac{a}{\ell},\frac{b}{\ell}[$ (où $\ell >0$).
\end{explanations}
\end{question}
}
